% Minimal retained preamble for lane-p2 cards taken from
% sources/hitting-sets-hypergraphs-2307.12154v3/source0.tex (arXiv:2307.12154v3).
%
% Copied verbatim from the author's own file:
%   line 3      : \documentclass[11pt,oneside]{amsart}
%   lines 43-49 : the author's macro definitions
%   lines 53-59 : the author's theorem declarations
% The base packages amsmath, amsfonts and amssymb are supplied here because amsart
% does not load amssymb by itself (LANE_BRIEF 18.1).
%
% Omitted author preamble lines and why (nothing was renamed or redefined):
%   line 10 (\usepackage[utf8]{inputenc}), line 11 (fontenc), line 12 (lmodern):
%       input encodings and font encodings fight the forced xelatex engine;
%   line 13 (geometry), line 34 (\baselinestretch), line 29 (parskip): pure layout,
%       and the assembler supplies the page setup;
%   lines 14-24 and 28 and 32 (array, fancyhdr, graphicx, float, caption, pgf, tikz,
%       tikzlings-addons, tikzlings-penguins, subcaption, adjustbox, verbatim): none of
%       the retained spans draws or floats anything, and the tikzlings package is not
%       part of the host TeX Live;
%   line 62 (\freefootnote): the retained spans contain no footnote.
\documentclass[11pt,oneside]{amsart}
\usepackage{amsmath}
\usepackage{amsfonts}
\usepackage{amssymb}

\newcommand{\R}{\mathbb{R}}
\newcommand{\N}{\mathbb{N}}
\newcommand{\HH}{\mathcal{H}}
\newcommand{\RR}{\mathcal{R}}
\newcommand{\A}{\mathcal{A}}
\newcommand{\B}{\mathcal{B}}
\newcommand{\est}{\emptyset}

\newtheorem{theorem}{Theorem}
\newtheorem{proposition}{Proposition}
\newtheorem{corollary}{Corollary}
\newtheorem{lemma}{Lemma}
\newtheorem{definition}{Definition}
\newtheorem{jel}{Notation}
\newtheorem{conj}{Conjecture}
\usepackage{graphicx}
\usepackage{verbatim}
\usepackage{pgf,tikz}
\usetikzlibrary{arrows}
\title{The polychromatic threshold of unions of a horizontal and a vertical strip lies between three quarters and four k}
\author{Balázs Bursics, Bence Csonka, Luca Szepessy}
\date{arxiv:2307.12154v3; CC BY 4.0}
\begin{document}
\maketitle
\noindent\textit{Editorial note.} The counted unit is the article's own Theorem kereszt at source lines 243-245 together with the proof environment at source lines 595-676 whose header names it: Proof of Theorem kereszt. The pairing is non-adjacent, because the article collects the introduction's proofs in a later section.\par\medskip
\noindent\textit{Editorial note.} The proof displays its construction as a tikz picture inside a figure environment whose caption names the theorem and whose label the proof references; the whole environment is inside the retained span, so the reference resolves within the material, and the retained preamble therefore carries the author's tikz load with the arrows library. No png asset is needed for this card.\par\medskip
\noindent\textit{Editorial note.} The proof cites the article on colorful strips once, in its upper-bound sentence. That entry is supplied as a task-directory fragment holding it exactly as a bibtex run on the host had formatted it from the author's own ref.bib; the bibtex backend itself cannot be used because the harness runs both xelatex passes before bibtex, so a generated bbl is never read back.\par\medskip
\noindent\textit{Editorial note.} The retained context at source lines 120, 128-129, 203-204 and 234-235 is the article's own definition of a polychromatic k-colouring, of the parameter m\_H(k), of the strip family with its dual and of the family of unions of two perpendicular strips.\par\medskip
\noindent\textit{Editorial note.} The retained preamble's remaining lines are the author's class line at source line 3, his macro definitions at source lines 43-49 and his theorem declarations at source lines 53-59, with amsmath, amsfonts and amssymb loaded explicitly for amsart.\par\medskip
\noindent\textit{Allowed standard prerequisites (editorial).}\begin{itemize}
\item The family of unions of a horizontal and a vertical strip. The family A plus consists of the hypergraphs on a finite vertex set in the plane whose edges are the intersections of the vertex set with the union of one horizontal and one vertical axis-parallel strip; the author's own definition is retained at source lines 234-235.. It is the family whose threshold the counted theorem bounds from both sides.
\item Polychromatic colourings and the parameter m\_H(k). A polychromatic k-colouring colours the vertices so that every edge meets all k colour classes, and m\_H(k) is the smallest threshold m such that every member of the family admits such a colouring after deleting the edges of size below m; retained at source lines 120 and 128-129.. The counted theorem brackets this parameter for the family of unions of two perpendicular strips.
\item The cited upper bound for the family of axis-parallel strips. The proof of the upper bound appeals to a theorem of Aloupis, Cardinal, Collette, Imahori, Korman, Langerman, Schwartz, Smorodinsky and Taslakian stating that any point set in the plane can be coloured so that all axis-parallel strips of size at least two k minus one are polychromatic; that entry is cited in the proof and is supplied with this card as a fragment of the bibtex-formatted entry.. It is the only external result used, and it is what makes the upper bound immediate: an edge of size at least four k minus three contains a strip of size at least two k minus one.
\end{itemize}\medskip
A {\it polychromatic k-coloring} of a hypergraph $H$ is a $k$-coloring of its vertex set such that every hyperedge contains a vertex of all $k$ color classes. Observe that a polychromatic 2-coloring is the same as the usual proper 2-coloring of hypergraphs, where we require that no edge is monochromatic. By merging $j$ color classes of a polychromatic $k$-coloring we get a polychromatic $(k-j+1)$-coloring, so the condition of monochromatic $k$-colorability becomes stricter as $k$ increases. A trivial necessary condition for the existence of a polychromatic $k$-coloring is that all edges of $H$ must be of size at least $k$.

\begin{definition} Let $\mathcal{H}$ be a hypergraph family. Denote by $m_\mathcal{H}(k)$ the smallest positive integer $m$ such that for every $H\in\mathcal{H}$ there exists a polychromatic $k$-coloring of the hypergraph $H_{\ge m}$. If there is no such $m$, set $m_{\mathcal{H}}(k)=\infty$.
\end{definition}

\begin{definition} Denote by $\A$ the hypergraph family which consists of those hypergraphs $H=(V,E)$ for which $V\subset \R^2$ is finite, and the edges are induced by axis-parallel strips: every $e\in E$ edge can be presented in the form $V\cap\{(x,y)\in\R^2: x_0<x<x_1\}$ or $V\cap\{(x,y)\in\R^2: y_0<y<y_1\}$ for some $x_0, x_1, y_0, y_1 \in\R$. Denote by $\A^*$ its dual: the family of such hypergraphs that the vertices are axis-parallel strips, and the edges are sets consisting of the strips containing a given point $(x,y)$. 
\end{definition}

\begin{definition} Denote by $\A^+$ the hypergraph family which consists of such hypergraphs $H=(V,E)$ for which $V\subset \R^2$ is finite, and the edges are induced by the union of a horizontal and a vertical axis-parallel strip: every $e\in E$ edge can be presented in the form $V\cap\big(\{(x,y)\in\R^2: y_0<y<y_1\}\cup \{(x,y)\in\R^2: x_0<x<x_1\}\big)$ for some $x_0,x_1,y_0,y_1\in\R$. 
\end{definition}

\begin{theorem}\label{kereszt}
$$3\bigg\lceil \frac{3}{4}k\bigg\rceil -2\le m_{\A^+}(k)\le 4k-3.$$
\end{theorem}

\begin{proof}[Proof of Theorem~\ref{kereszt}]

For proving the upper bound, color an arbitrary point set in the plane in such a way that all axis-parallel strips of size at least $2k-1$ are polychromatic, this is possible by Theorem 1. in \cite{apstrip}. Then any edge of size at least $4k-3$ contains a horizontal or vertical axis-parallel strip of size at least $2k-1$ and thus polychromatic, implying that the coloring is polychromatic.

For the construction to prove the lower bound, take 8 sets $H_1, H_2, \ldots, H_8$, each consisting of $\lceil \frac{3}{4}k\rceil-1$ points in such an arrangement so that for $(x_i,y_i)\in H_i$ (where $1\le i\le8$) we have 

\begin{center}
    $x_1<x_2<x_3<x_4<x_5<x_6<x_7<x_8$

$y_7<y_5<y_8<y_6<y_3<y_1<y_4<y_2,$  
\end{center}

\noindent see Figure \ref{keresztabra}.

\begin{figure}[h!]
\centering
\begin{tikzpicture}[line cap=round,line join=round,>=triangle 45,x=0.8cm,y=0.8cm]
\clip(-0.5,-3.5) rectangle (9.5,4.5);
\draw(3,1) circle (0.24cm);
\fill [color=black] (3.1,1.1) circle (0.5pt);
\fill [color=black] (3,0.9) circle (0.5pt);
\fill [color=black] (2.9,1) circle (0.5pt);
\fill [color=black] (2.95,1.05) circle (0.5pt);
\draw(2,4) circle (0.24cm);
\fill [color=black] (2.1,4.1) circle (0.5pt);
\fill [color=black] (2,3.9) circle (0.5pt);
\fill [color=black] (1.9,4) circle (0.5pt);
\fill [color=black] (1.95,4.05) circle (0.5pt);
\draw(4,3) circle (0.24cm);
\fill [color=black] (4.1,3.1) circle (0.5pt);
\fill [color=black] (4,2.9) circle (0.5pt);
\fill [color=black] (3.9,3) circle (0.5pt);
\fill [color=black] (3.95,3.05) circle (0.5pt);
\draw(5,-2) circle (0.24cm);
\fill [color=black] (5.1,-2.1) circle (0.5pt);
\fill [color=black] (5,-1.9) circle (0.5pt);
\fill [color=black] (4.9,-2) circle (0.5pt);
\fill [color=black] (4.95,-2.05) circle (0.5pt);
\draw(1,2) circle (0.24cm);
\fill [color=black] (1.1,2.1) circle (0.5pt);
\fill [color=black] (1,1.9) circle (0.5pt);
\fill [color=black] (0.9,2) circle (0.5pt);
\fill [color=black] (0.95,2.05) circle (0.5pt);
\draw(6,0) circle (0.24cm);
\fill [color=black] (6.1,0.1) circle (0.5pt);
\fill [color=black] (6,-0.1) circle (0.5pt);
\fill [color=black] (5.9,0) circle (0.5pt);
\fill [color=black] (5.95,0.05) circle (0.5pt);
\draw(7,-3) circle (0.24cm);
\fill [color=black] (7.1,-3.1) circle (0.5pt);
\fill [color=black] (7,-2.9) circle (0.5pt);
\fill [color=black] (6.9,-3) circle (0.5pt);
\fill [color=black] (6.95,-3.05) circle (0.5pt);
\draw(8,-1) circle (0.24cm);
\fill [color=black] (8.1,-1.1) circle (0.5pt);
\fill [color=black] (8,-0.9) circle (0.5pt);
\fill [color=black] (7.9,-1) circle (0.5pt);
\fill [color=black] (7.95,-1.05) circle (0.5pt);

\draw (1.5,2.5) node [anchor=north west][inner sep=0.75pt]    {$H_1$};

\draw (2.5,4.5) node [anchor=north west][inner sep=0.75pt]    {$H_2$};

\draw (3.5,1.5) node [anchor=north west][inner sep=0.75pt]    {$H_3$};

\draw (4.5,3.5) node [anchor=north west][inner sep=0.75pt]    {$H_4$};

\draw (5.5,-1.5) node [anchor=north west][inner sep=0.75pt]    {$H_5$};

\draw (6.5,0.5) node [anchor=north west][inner sep=0.75pt]    {$H_6$};

\draw (7.5,-2.5) node [anchor=north west][inner sep=0.75pt]    {$H_7$};

\draw (8.5,-0.5) node [anchor=north west][inner sep=0.75pt]    {$H_8$};

\end{tikzpicture}
\caption{Construction for Theorem~\ref{kereszt}.}
\label{keresztabra}
\end{figure}

For an arbitrary coloring of the vertices each set out of the eight has strictly more than $\frac{k}{4}$ colors which does not appear in them, so by the pigeonhole principle we have 3 sets with a common missing color. Observe that any three sets can be separated from the others with the union of a horizontal and a vertical strip, which implies $m(k)>3(\lceil \frac{3}{4}k\rceil-1)$.
\end{proof}

\begin{thebibliography}{99}
\bibitem{apstrip}
G.~Aloupis, J.~Cardinal, S.~Collette, S.~Imahori, M.~Korman, S.~Langerman,
  O.~Schwartz, S.~Smorodinsky, and P.~Taslakian.
\newblock Colorful strips.
\newblock {\em Graphs and combinatorics}, 27:327--339, 2011.
\end{thebibliography}
\end{document}
